\section{Xây dựng Edge Agent}

\subsection{Lắp ráp và cấu hình}

\texttt{AgentContainer} lắp ráp Agent từ cấu hình \texttt{AgentSettings} (tiền tố
biến môi trường \texttt{EDGE\_}). Các lựa chọn quan trọng: \texttt{runtime\_driver}
(\texttt{docker} hoặc \texttt{fake} để chạy không cần Docker), \texttt{runtime\_network},
\texttt{image\_pull\_policy}, \texttt{agent\_url\_for\_runtimes} (địa chỉ để container
runtime gọi ngược về Agent, mặc định \texttt{host.docker.internal:8081}),
\texttt{state\_dir} lưu desired state và cache model, cùng ba chu kỳ state (10~s),
telemetry (5~s), reconcile (15~s). Ứng dụng FastAPI của Agent được tạo bởi
\texttt{AgentAppFactory}: \texttt{lifespan} khởi động AgentService và dừng an
toàn khi tắt; health check báo trạng thái kết nối broker.

\subsection{Giao tiếp MQTT}

\texttt{EdgePublisher} đóng gói việc publish NodeState (retained), telemetry và
event lên topic đúng định dạng, dùng \texttt{EdgeTopics} chung. Vì broker có thể
giao lại thông điệp, tính đúng đắn không dựa vào QoS mà dựa vào idempotency ở
tầng ứng dụng: chống trùng \texttt{command\_id} và thao tác runtime idempotent.
Khi Agent tạm mất kết nối, AI Runtime vẫn chạy bình thường vì không phụ thuộc
Cloud; vòng reconcile tiếp tục bảo đảm container sống; sau khi kết nối lại, bản
tin NodeState kế tiếp đồng bộ lại trạng thái trên Cloud.

\subsection{Điều khiển container}

\texttt{DockerRuntimeDriver} hiện thực port \texttt{RuntimeDriver} bằng cách gọi
Docker CLI qua \texttt{asyncio.create\_subprocess\_exec}, hoạt động được cả khi
Agent chạy trong container có gắn socket. Container được đặt tên
\texttt{edgeops-<tên>-<8 ký tự đầu runtime\_id>} và gắn nhãn
\texttt{edgeops.runtime\_id} để tìm lại sau khi Agent khởi động lại. Hàm
\texttt{inspect} đọc trạng thái chạy, mã thoát, số lần restart, thời điểm khởi
động và lỗi từ \texttt{docker inspect}. \texttt{FakeRuntimeDriver} mô phỏng cùng
giao diện trong bộ nhớ, cho phép kiểm thử kịch bản container bị crash.

\subsection{Tải model}

\texttt{HttpArtifactFetcher} tải model theo luồng (stream) với thời gian chờ 300~s,
ghi vào tệp \texttt{.part} rồi \texttt{os.replace} sang tên đích để một lần tải dở
dang không bao giờ bị dùng nhầm. Thư mục cache được phân tách theo 16 ký tự đầu
của SHA-256 URL. Lược đồ URI không hỗ trợ (ví dụ \texttt{s3://}) bị từ chối với
thông báo yêu cầu dùng URL HTTPS có chữ ký. Việc tải được thực hiện trong bước
CREATING của runtime nên lỗi tải được báo về Cloud như một deployment thất bại.

\subsection{Thu thập telemetry}

\texttt{PsutilSystemCollector} thu CPU, tải trung bình 1/5 phút, iowait, bộ nhớ,
swap, nhiệt độ (giá trị lớn nhất của các cảm biến hệ điều hành cung cấp), dung
lượng đĩa và inode, uptime, số tiến trình. Các chỉ số dạng tốc độ --- số lần chuyển
ngữ cảnh mỗi giây, tốc độ tăng bộ nhớ (MB/phút), thông lượng đọc/ghi đĩa và
mạng (MB/s), tốc độ tăng nhiệt (°C/phút) --- được tính từ chênh lệch với lần đo
trước.
\texttt{NvidiaSmiGpuCollector} truy vấn \texttt{nvidia-smi} để lấy mức sử dụng,
bộ nhớ, xung nhịp, công suất và nhiệt độ GPU; các giá trị \texttt{[N/A]} được
chuyển thành rỗng. \texttt{PsutilNodeInfoProvider} cung cấp hostname, IP, số nhân,
dung lượng bộ nhớ và tên GPU cho NodeState.
